\chapter{Introduction}
\label{chap:introduction}

\section{Industrial Context and Problem Statement}
Aviation and industrial turbofan engines are complex aeromechanical systems that operate under extreme thermo-mechanical stresses, severe vibrational loads, and dynamically varying flight envelopes \cite{saxena2008damage}. In contemporary fleet operations, ensuring operational reliability while managing lifecycle maintenance expenditures represents a paramount economic and safety challenge. Unplanned in-flight engine shutdowns, emergency diversions, and catastrophic component structural failures impose severe safety risks and immense financial costs, frequently exceeding scheduled maintenance overhauls by orders of magnitude \cite{chen2022data}.

Historically, commercial aviation maintenance relied predominantly on conservative time-based preventative policies, replacing subassemblies at predetermined flight-hour or cycle milestones regardless of actual physical health \cite{shehadeh2024evaluating}. While safe, this schedule-driven paradigm leads to premature replacement of functional hardware, wasting significant operational life. 

With the emergence of Internet of Things (IoT) sensor telemetry and high-capacity flight data recorders, Condition-Based Maintenance (CBM) and Predictive Maintenance (PdM) paradigms have become foundational to next-generation aviation operations. These paradigms seek to predict the Remaining Useful Life (RUL) of critical components dynamically from multivariate sensor streams, enabling opportunistic scheduling before physical failure occurs \cite{huang2021survey}.

\section{The Prognostic-to-Decision Gap}
Despite significant advances in deep learning for degradation modeling, an acute methodological gap persists between cycle-level machine learning accuracy and downstream operational utility:
\begin{enumerate}
    \item \textbf{Cycle-Level Error Optimization}: Point regression architectures—such as Long Short-Term Memory (LSTM) networks, Temporal Convolutional Networks (TCNs), and Gradient Boosted Regressors (GBRs)—are conventionally trained to minimize symmetric cycle-level loss functions such as Mean Squared Error (MSE) or Mean Absolute Error (MAE) \cite{saxena2008damage, diao2026turbofan}. These metrics treat underestimating and overestimating RUL symmetrically.
    \item \textbf{Asymmetric Cost Realities}: In real-world aviation logistics, overestimating RUL by 10 cycles can allow an engine to operate beyond its physical failure boundary, causing catastrophic damage. In contrast, underestimating RUL by 10 cycles merely results in premature maintenance and a minor loss of usable component life \cite{saxena2008damage, robinson2026riskaware}.
    \item \textbf{Uncalibrated Uncertainty}: Predictive uncertainty methods differ in their assumptions and calibration behavior. Their maintenance value must be evaluated through interval reliability and downstream decision outcomes under a common protocol.
\end{enumerate}

\section{The IntelliTwin Approach}
To resolve this gap, the \textbf{IntelliTwin} research prototype integrates data-driven base degradation predictors with hierarchical, decision-aware conformal prediction theory.

IntelliTwin uses an engine-weighted conformal construction motivated by hierarchical conformal prediction. Formal coverage results require the hierarchical sampling and exchangeability assumptions stated in the underlying theory. Equal engine weighting does not establish these assumptions for ordered, lifetime-stratified degradation trajectories. This study therefore evaluates marginal and subgroup coverage empirically. It does not claim guaranteed coverage for each engine, lifecycle stage, operating regime, complete trajectory, or adaptively selected maintenance time. Nominal 90% coverage is a target, not a guaranteed realized coverage proportion. Maintenance outcomes are evaluated separately under the frozen simulator.

Furthermore, the framework introduces two critical adaptations:
\begin{itemize}
    \item \textbf{Regime-Stratified Calibration}: Conformal calibration is conditioned on operational flight regime clusters, evaluating whether localized calibration improves conditional coverage under stressed operating conditions.
    \item \textbf{Decision-Aware Asymmetry}: Nonconformity scores are scaled via an asymmetric parameter ($\beta = -\ln 2$), tilting prediction intervals to guard against hazardous RUL overestimation.
\end{itemize}
The system is evaluated as an offline, trace-driven prognostic decision-support research prototype; physical engine integration and operational airline validation have not been demonstrated, and maintenance costs are evaluated in Simulated Cost Units (SCU) reflecting synthetic benchmark preferences rather than validated airline economics.

\section{Research Objectives and Scope}
This project report / thesis draft formalizes four prespecified research questions version-controlled prior to held-out evaluation:
\begin{itemize}
    \item \textbf{RQ1 (Sequential Utility)}: \textit{Can a decision-aware, condition-adaptive conformal RUL calibration method improve sequential maintenance utility while preserving reliable prediction-interval coverage?}
    \item \textbf{RQ2 (Subgroup Reliability)}: \textit{Can the proposed method improve reliability across operating regimes and degradation conditions without excessively widening prediction intervals?}
    \item \textbf{RQ3 (Attribution vs. Matched Controls)}: \textit{Are the maintenance benefits of the proposed method attributable to decision-aware uncertainty calibration rather than simply more conservative prediction intervals or earlier maintenance?}
    \item \textbf{RQ4 (Robustness \& Generalization)}: \textit{Does the proposed method retain its calibration and maintenance-decision benefits across operating conditions, base RUL predictors, datasets, and repeated experimental runs?}
\end{itemize}

\section{Thesis Organization}
The remainder of this report/thesis draft is organized as follows:
\begin{itemize}
    \item Chapter~\ref{chap:literature_review} provides a comprehensive review of turbofan prognostics, conformal prediction, and maintenance decision theory.
    \item Chapter~\ref{chap:problem_hypotheses} details the mathematical problem formulation and authoritative hypothesis framework.
    \item Chapter~\ref{chap:methodology} presents the complete methodological pipeline from causal windowing to asymmetric calibration.
    \item Chapter~\ref{chap:implementation} outlines the software architecture, engine-level partitioning, and provenance tracking.
    \item Chapter~\ref{chap:experimental_design} formalizes the confirmatory experimental setup, baselines, and bootstrap inference.
    \item Chapter~\ref{chap:results_discussion} presents the status of experimental verification and discussion of expected outcomes.
    \item Chapter~\ref{chap:limitations_conclusion} discusses threats to validity, research boundaries, and conclusions.
    \item Appendix~\ref{chap:reproducibility} provides the complete reproducibility protocol and execution commands.
\end{itemize}
